% !TEX TS-program = latex
% -*- mode : latex; coding: utf-8 -*-

%% 재단 맡길 때에는 showtrims을 켜서 a4로 출력 후 잘라내면 됨.
%\documentclass[11pt,a4paper,oneside,openany,reqno,chapter,nokorean,showtrims]{oblivoir}
\documentclass[11pt,a4paper,oneside,openany,reqno,chapter,nokorean]{oblivoir}

\usepackage[utf8]{inputenc}
\usepackage[dbl4x6]{fapapersize}

\chapterstyle{default}
\usepackage{tocloft}% http://ctan.org/pkg/tocloft
\def\cftchapteraftersnum{.}
\renewcommand{\cftchapterdotsep}{\cftdotsep}% Chapters should use dots in ToC


\let\interHchar\undefined
\let\interhword\undefined
\usepackage[default]{dhucs-interword}
\usepackage{dhucs-gremph}
\usepackage{dhucs-enumerate}
%\usepackage{enumerate}
\usepackage{ifpdf}
\usepackage{url}
\usepackage{rotating}



\ifpdf
\usepackage[pdftex]{hyperref} 
\hypersetup{
colorlinks = true
} 
\fi

\usepackage{amsmath,amsthm,amssymb,amsxtra}
\usepackage{graphicx}
\newsubfloat{figure}
\usepackage{epsfig}
\let\footruleskip\undefined
\usepackage{fancyhdr}
\usepackage{environ}
\usepackage{siunitx}
\ifpdf
% enable hyper links
\hypersetup{
 pdftitle={Thesis},
 pdfauthor={Geonmo Ryu},
 pdfkeywords={j/psi,charmed meson},
 citecolor=black,
 %citecolor=darkblue,
 filecolor=black,
 linkcolor=black,
 %inkcolor=darkblue,
 urlcolor=blue,
 colorlinks=true}
\fi
\usepackage{uosphysicsEnglish} %서울시립대양식

\setlength\parindent{2pc} % 들여쓰기 1글자

\newtheorem{theorem}{Theorem}[chapter] % Theorem 1.1, Theorem 3.4
\newtheorem{corollary}[theorem]{Corollary} % Corollary는 독자번호 없음.
\newtheorem{lemma}[theorem]{Lemma}
\newtheorem{proposition}[theorem]{Proposition}
\newtheorem{question}{Question} % 독자적인 번호를 갖는다. Question 1, Question 2
\newtheorem{step}{Step}

% 번호를 붙이지 않는 `정리' 만들기
\newtheorem{lemmawon}{Lemma}
\newtheorem{claim}{Claim}

% 특별한 이름을 갖는, 번호가 붙지 않는 Theorem 만들기
\newtheorem{questionA}{Question~A}
\newtheorem{mainthmA}{Main Theorem~A}

% 본문의 글씨가 기울어지지 않는 환경
\theoremstyle{definition}
\newtheorem{definition}[theorem]{Definition}

% 화면에 나오는 이름이 이탤릭체인 환경
\theoremstyle{remark}
\newtheorem{remark}[theorem]{Remark}
\newtheorem{remarkof}{Remark}[theorem] % 어떤 Theorem의 아래에 Remark
\newtheorem{remarkwon}{Remark}

% 수식 번호 붙이기 (1.1), (3.3)
\numberwithin{equation}{chapter} 

%책갈피용
%\usepackage[bookmarks=true]{hyperref}
\usepackage{bookmark}


\setcounter{secnumdepth}{4}
\setcounter{tocdepth}{2}


%표 정렬용
\usepackage{multirow}
\renewcommand{\multirowsetup}{\centering}

%그림위 글자 기능(인준서 도장용)
\usepackage[abs]{overpic}
\setlength\unitlength{1mm}

\begin{document}
  
  %% 정보
  \EnglishTitle{Top quark mass measurement using charmed meson within b-jet}
  %% 본문이 영문일 경우, 국문 제목은 빈칸으로 둘 것
  %\KoreanTitle{b-제트 안의 참 쿼크 중간자를 이용한 탑 쿼크 질량 측정}
  \KoreanTitle{}
  \AuthorKoreanName{류 건 모}
  \AuthorEnglishName{Geonmo Ryu}
  \KoreanNameofDegree{물리학 박사}
  
  %% 해당 내용으로 채울 것. 석사학위의 경우 학과에 따라 "Science" 등 별도 용어 필요
  %% Mater of Physics 등도 심사를 맡아봐야 할듯.
  \EnglishNameofDegree{Doctor of Philosophy}
  \EnglishNameofPaper{Ph.D. Dissertation}

  %\EnglishNameofDegree{Master of ???}
  %\EnglishNameofPaper{Master Thesis}

  \GraduateDate{2017년 2월}
  \GraduateDateEnglish{February 2017}
  \SummittedDate{2016년 12월}
	\EnglishSummittedDate{December 2016}
  \EnglishDepartmentName{Department of Physics}
  \KoreanDepartmentName{물리학과}
  \RefereeDate{2016년 12월}
  \RefereeChief{Hyunsoo Min}
  \RefereeSecond{Inkyu Park}
  \RefereeThird{Jason Sanghun Lee}
	\RefereeFourth{Intae Yu}
	\RefereeFifth{Tae Jeong Kim}
  \AdvisorKoreanName{박 인 규}
  \AdvisorEnglishName{Inkyu Park}
  \KeyWordsEnglish{LHC, CMS, top quark, $J/\psi$, D0, D*, charmed meson}
  \KeyWordsKorean{LHC, CMS, 탑 쿼크, $J/\psi$ 중간자, D0, D*, 참 중간자}

  %% 앞표지
  \makefrontcover
  
  %% 
	\englishproposal
  
  %% 국문 제출 승인서. 국문 인준서의 경우 업데이트가 되지 않아서 추가 작업 필요.
  %\makeapproval

  %% 영문 제출 승인서
  \makeenglishapproval
  %% 국문초록
  \pagenumbering{gobble}
  \include{abs-english}
  \makefrontcover
  \cedp
  
  %% 본문 이전의 쪽들
  \frontmatter
  
  %% TOC/ LOF / LOT 끝의 * 지우면 안됨. ( 목차에서 명단 삭제 필요 )
  \pagestyle{plain}%
	%% 목차
	\clearpage
	\phantomsection 
    \tableofcontents*
    
	%% 그림목차
	\clearpage
	\phantomsection
    \listoffigures*
	%% 표목차
	\clearpage
	\phantomsection
	\listoftables*
  
  %%% 본문 시작
  \mainmatter
  \linespread{2}
  %\SetHangulspace{2.2}{2.0}

  %% 쪽번호가 아래 중앙에 붙는다.
  \pagestyle{fancy}%
  \lhead{\leftmark}%
  \rhead{}%
  \chead{}%
  \cfoot{\thepage}%
  \renewcommand{\headrulewidth}{0pt}
  
  % \part{1부 소개}
  \include{chapter1}
  \cedp
  \include{chapter2}
  \cedp
  \include{chapter3}
  \cedp
  \include{chapter4}
  \cedp
  \include{chapter5}
  \cedp
  \include{chapter6}
  \cedp
  \include{chapter7}
  \cedp
  \include{chapter8}
  \cedp
  \include{chapter9}
  \cedp

  \pagestyle{plain}%
  %% 참고문헌
	\clearpage
	\phantomsection	
  \include{biblo}
  \cedp
 
  %% 초록
  \include{abs-korean}	
  \cedp
  
  % 감사의 글
  \include{thanksto}
  \cedp
  

\end{document}